\documentclass{beamer}
\usepackage{graphicx}
\usepackage{xcolor}
\usepackage{tikz}

% 自定义颜色
\definecolor{purpleblue}{RGB}{80, 80, 180}
\definecolor{lightgray}{RGB}{230, 230, 230}
\definecolor{novelPurple}{RGB}{98, 54, 255} 
\definecolor{darkgray}{RGB}{50, 50, 50}

% 主题设置
\usetheme{Madrid}  
\usecolortheme{dolphin} 

% 设置字体（移除 fontspec，确保 pdfLaTeX 兼容）
\setbeamerfont{title}{size=\Large,series=\bfseries}
\setbeamerfont{frametitle}{size=\LARGE,series=\bfseries}
\setbeamercolor{frametitle}{fg=white, bg=novelPurple}
\setbeamercolor{title}{fg=white, bg=purpleblue}

% 封面信息
\title[SAT Hard Question]{\textbf{SAT Hard Question 5}}
\author{\textbf{Jack Zeng}}
\institute{\textcolor{white}{\textbf{Novel Prep}}}
\date{\today} 

\begin{document}

% --- Title Slide ---
\begin{frame}
    \titlepage
    \vfill
    \begin{flushright}
        \textcolor{purpleblue}{\Huge \textbf{Novel Prep}}
    \end{flushright}
\end{frame}


\begin{frame}
    \begin{tikzpicture}[remember picture, overlay]
        % 设置浅色渐变背景
        \shade[top color=softPurple, bottom color=softGray] 
            (current page.north west) rectangle (current page.south east);
        
        % 添加高斯模糊光晕
        \fill[white, opacity=0.3] (current page.center) circle (4cm);
        
        % 主标题
        \node[align=center] at (current page.center) {
            {\Huge\bfseries\textcolor{purpleblue}{Novel Prep}} \\[0.5cm]
            {\Large\itshape\textcolor{lightText}{Excellence in SAT Prep}}
        };

        % 添加柔和光点（点缀）
        \foreach \x in {1,2,3,4,5} {
            \fill[purpleblue, opacity=0.2] (rand*10-5, rand*7-3.5) circle (0.1+\x*0.05);
        }

        ;
    \end{tikzpicture}
\end{frame}


% Question 19
\begin{frame}{Question: Expression}
    
The given expression can be rewritten as \(\frac{a}{4}x^2 + \frac{b}{4}x + \frac{c}{4}\) where \(a\), \(b\), and \(c\) are constants. What is the value of \(a+b+c\)? Given the expression \(-16(5x-3)^2 + 4(5x-2)^2\).
\begin{itemize}
    \item[A] -112
    \item[B] -56
    \item[C] -28
    \item[D] -7
\end{itemize}
\end{frame}

\begin{frame}{Answer}
    

\textbf{Correct Answer: A.} \\
The coefficients \(a\), \(b\), and \(c\) are derived from the given quadratic expression by expanding and combining like terms. Calculating \(a+b+c\) yields \(-112\).
\end{frame}

\begin{frame}{Question 20}
In the given function, $b$ is a positive constant. For every increase in the value of $x$ by 1, the value of $f(x)$ increases by $c\%$ where $0 < c < 100$. Which expression gives the value of $c$ in terms of $b$?
\[
f(x) = 66(b)^x
\]
\begin{itemize}
    \item[A.] $1 + \frac{b}{100}$
    \item[B.] $b + 100$
    \item[C.] $100(b + 1)$
    \item[D.] $100(b - 1)$
\end{itemize}
\end{frame}

\begin{frame}{Answer to Question 20}
\textbf{Correct Answer: D.} \\
Given the exponential function $f(x)$, the percentage change for an increment of 1 in $x$ corresponds to $100(b - 1)\%$.
\end{frame}

\begin{frame}{Question 11}
If $\sin x^\circ = a$, which of the following must be true for all values of $x$?
\begin{itemize}
    \item[A.] $\cos x^\circ = a$
    \item[B.] $\sin(90^\circ - x^\circ) = a$
    \item[C.] $\sin (x^2)^\circ = a^2$
    \item[D.] $\cos(90^\circ - x^\circ) = a$
\end{itemize}
\end{frame}

\begin{frame}{Answer to Question 11}
\textbf{Correct Answer: D.} \\
This is the complementary angle identity: $\sin(\theta) = \cos(90^\circ - \theta)$.
\end{frame}

\begin{frame}{Question 20}
The function $g$ is a quadratic function in the $xy$-plane. The graph of $y = g(x)$ has a vertex at $(-1, -4)$ and passes through the points $(-2, -43)$ and $(1, -160)$.  
What is the value of $g(0) - g(2)$?
\begin{itemize}
    \item[A.] 312
    \item[B.] 355
    \item[C.] -117
    \item[D.] -203
\end{itemize}
\end{frame}

\begin{frame}{Answer to Question 20}
\textbf{Correct Answer: A.}
\end{frame}

\begin{frame}{Question 22}
The mass of object A is $324\%$ of the mass of object B and the mass of object A is $7.2\%$ of the mass of object C.  
If the mass of object C is $p\%$ of the mass of object B, what is the value of $p\%$?
\begin{itemize}
    \item[A.] $4.5\%$
    \item[B.] $4.5$
    \item[C.] $45$
    \item[D.] $450$
\end{itemize}
\end{frame}

\begin{frame}{Answer to Question 22}
\textbf{Correct Answer: C.}
\end{frame}

\begin{frame}{Question 22}
In the given equation, $c$ is a constant. If the equation has exactly one solution, what is the value of $c$?
\[
\frac{1}{cx} = \frac{x}{76} + \frac{1}{c}
\]

\end{frame}

\begin{frame}{Answer to Question 22}
\textbf{Correct Answer: -19.}
\end{frame}

\begin{frame}{Question 22}
A parallel electric circuit contains four resistors with resistances of $p$ ohms, $q$ ohms, $s$ ohms, and 13 ohms. The given equation relates the resistance $R$, in ohms, of the circuit, to the resistances of the four individual resistors it contains. Which equation correctly expresses $R$ in terms of $p$, $q$, and $s$?
\[
\frac{1}{R} = \frac{1}{p} + \frac{1}{q} + \frac{1}{s} + \frac{1}{13}
\]
\begin{itemize}
    \item[A.] $R = \frac{13pqs}{p+q+s+13}$
    \item[B.] $R = \frac{13pqs}{13qs + 13ps + 13pq + pqs}$
    \item[C.] $R = \frac{13qs + 13ps + 13pq + pqs}{13pqs}$
    \item[D.] $R = p + q + s + 13$
\end{itemize}
\end{frame}

\begin{frame}{Answer to Question 22}
\textbf{Correct Answer: B.} \\
This follows from taking the reciprocal of the sum of reciprocals of the resistors in a parallel circuit.
\end{frame}

\begin{frame}{Question 21}
In the given equation, $b$ is an integer. The equation has no real solutions.  
What is the least possible value of $b$?
\[
-x^2 + bx - 169 = 0
\]
\begin{itemize}
    \item[A.] 26
    \item[B.] 25
    \item[C.] -26
    \item[D.] -25
\end{itemize}
\end{frame}

\begin{frame}{Answer to Question 21}
\textbf{Correct Answer: D.} \\
The discriminant $b^2 - 4ac$ must be less than 0 for no real solutions. Solve for minimum $b$ such that $b^2 < 4(-1)(-169)$.
\end{frame}

\begin{frame}{Question 19}
The exponential function $f$ is defined by $f(x) = ab^x$, where $a$ and $b$ are positive constants.  
If $f(n - 1) = f(n) + \frac{82}{100}f(n - 1)$, where $n$ is a constant, what is the value of $b$?
\begin{itemize}
    \item[A.] 1
    \item[B.] 0.82
    \item[C.] 0.18
    \item[D.] 10
\end{itemize}
\end{frame}

\begin{frame}{Answer to Question 19}
\textbf{Correct Answer: C.} \\
Use the equation recursively to isolate $b$ and simplify to find the correct constant multiplier.
\end{frame}

\begin{frame}{Question 13}
A rectangle is inscribed in a circle such that the length of the diagonal of the rectangle is twice the length of its shortest side.  
The circumference of the circle is $114\pi$ units.  
What is the area, in square units, of the rectangle?
\begin{itemize}
    \item[A.] $57\sqrt{2}$
    \item[B.] $57\sqrt{3}$
    \item[C.] $3,\!249\sqrt{2}$
    \item[D.] $3,\!249\sqrt{3}$
\end{itemize}
\end{frame}

\begin{frame}{Answer to Question 13}
\textbf{Correct Answer: D.} \\
Use the relationship between diagonal and side length, then relate it to the radius (half the diagonal) from the circumference to solve for area.
\end{frame}

\newpage

\begin{frame}{Question 15}
In triangle $PQR$, the measure of angle $P$ is $(3x + 5)^\circ$, the measure of angle $Q$ is $(2x + 9)^\circ$, and the measure of angle $R$ is $(4y + 5)^\circ$.  
If side $QR$ is extended through point $R$ to point $S$, and the measure of angle $PRS$ is $(x + y)^\circ$, what is the value of $x + y$?

\begin{itemize}
    \item[A.] 39
    \item[B.] 29
    \item[C.] 141
    \item[D.] 146
\end{itemize}
\end{frame}

\begin{frame}{Answer to Question 15}
\textbf{Correct Answer: A.}  
\end{frame}


\begin{frame}{Question 11}
The manager of a gym selected a sample of 139 members at random to estimate the percentage of the gym's members that would continue to pay for a membership if the price increased.  
From the survey, the manager estimates that $82\%$ of the gym's members would continue to pay, with a margin of error of $6.39\%$.  
If the survey is repeated with a random sample of 278 members and the results are calculated the same way, what is the most likely effect of using the larger sample?

\begin{itemize}
    \item[A.] The margin of error will be lower.
    \item[B.] The margin of error will be higher.
    \item[C.] The estimate will be lower.
    \item[D.] The estimate will be higher.
\end{itemize}
\end{frame}


\begin{frame}{Question 18}
Given the expression $\sqrt[5]{119n \left(\sqrt[6]{119n}\right)^2}$, for what value of $x$ is this equivalent to $(119n)^{30x}$, where $n > 1$?

\begin{itemize}
    \item[A.] 225
    \item[B.] $\frac{8}{225}$
    \item[C.] $\frac{4}{225}$
    \item[D.] $\frac{2}{225}$
\end{itemize}
\end{frame}

\begin{frame}{Answer to Question 18}
\textbf{Correct Answer: D.}
\end{frame}


\begin{frame}{Question 21}
The mass of object A is $483\%$ of the mass of object B, and the mass of object A is $0.069\%$ of the mass of object C.  
If the mass of object C is $p\%$ of the mass of object B, what is the value of $\frac{p}{1,000}$?

\begin{itemize}
    \item[A.] $0.7\%$
    \item[B.] 70
    \item[C.] $70\%$
    \item[D.] 700
\end{itemize}
\end{frame}

\begin{frame}{Answer to Question 21}
\textbf{Correct Answer: D.}
\end{frame}


\begin{frame}{Question 22}
For triangle $PQR$ and triangle $STU$, angles $P$ and $S$ each measure $30^\circ$, with $QR = 16$ and $TU = 12$.  
Which additional piece(s) of information is(are) sufficient to prove triangle $PQR$ is similar to triangle $STU$?

\begin{itemize}
    \item[I.] $Q = 40^\circ$, $U = 110^\circ$
    \item[II.] $\frac{PR}{SU} = \frac{4}{3}$
\end{itemize}

\begin{itemize}
    \item[A.] I is sufficient, but II is not.
    \item[B.] II is sufficient, but I is not.
    \item[C.] I is sufficient, and II is sufficient.
    \item[D.] Neither I nor II is sufficient.
\end{itemize}
\end{frame}






\end{document}
